\section{Ameaças à Validade}
\label{sec:ameacas}

Há ameaças importantes à validade dos resultados. As carteiras foram
obtidas de fontes públicas e podem não representar a distribuição
completa de carteiras reais de investidores; os perfis do workload,
além disso, não estão perfeitamente balanceados entre conservador,
moderado e agressivo. A validação estrutural não utiliza dados externos,
o que favorece diagnósticos inconclusivos nos sistemas que tratam
ausência de evidência com mais rigor, e deve ser lida como avaliação
estrutural, não como comparação principal de qualidade financeira.

A comparação enriquecida, por sua vez, foi feita em 20 carteiras. Ela é a
base da comparação entre técnicas, mas ainda é uma amostra limitada para
afirmar superioridade estatística de uma arquitetura sobre outra. Além
disso, parte dessa execução precisou ser realizada com menor concorrência
por instabilidade temporária do provedor, o que afeta diretamente as
métricas de latência. Os modelos usados via provedor externo podem variar
de comportamento ao longo do tempo, mesmo com prompts e entradas
congeladas. A métrica de acurácia proxy por consenso também pode favorecer
abordagens que se comportam de forma semelhante à maioria, sem garantir que
a maioria esteja correta financeiramente. E esta versão ainda não inclui
avaliação humana, LLM-as-judge calibrado, verificação detalhada de erro
factual nem comparação contra retorno futuro das carteiras.

Para reduzir esses riscos, o pipeline registra todas as chamadas e
resultados em SQLite, usa entradas congeladas, compara múltiplos
baselines, reporta erros de parsing e separa métricas operacionais de
métricas de qualidade. Mesmo assim, as conclusões aqui apresentadas
devem ser lidas como evidência inicial sobre a arquitetura e seu
comportamento, não como validação definitiva de desempenho financeiro.
